\chapter{Linear Dynamics}\label{linear:ch:dynamics}

\begin{center}\itshape
The motion of a machine with its governor consists in general of a uniform motion, combined with a disturbance which may be
expressed as the sum of several component motions. These components may be of four different kinds.
\upshape(James Clerk Maxwell, \emph{On Governors}, 1868)
\end{center}

In the 1780s James Watt fitted his steam engines with a \emph{centrifugal governor}: two heavy balls on hinged arms, spun by the
engine, that rise when it runs too fast and close the steam valve through a lever. It was the first widely used feedback device,
and by the 1860s there were tens of thousands of them, many of which misbehaved. Engines with governors would ``hunt'', speeding up
and slowing down in growing oscillations, and engineers tuned them by trial and error. In 1868 James Clerk Maxwell, between his
papers on electromagnetism, wrote a short note \emph{On Governors} for the Royal Society. He linearised the equations of motion
near the steady state, wrote the disturbance as a sum of exponentials $e^{\lambda t}$, and observed that the governor is stable
exactly when every $\lambda$ has negative real part. For his third-order model he found the condition explicitly; for higher
orders he could not, and he posed the problem publicly. Edward Routh solved it in 1877 and won the Adams Prize for it; Adolf
Hurwitz found the same criterion independently in 1895, asked by an engineer designing turbines in Switzerland.

Maxwell's observation is the heart of this chapter: the long-term behaviour of a linear system is written in the eigenvalues of its
matrix. For iterations $x_{k+1}=Ax_k$ the test is $|\lambda|<1$, for differential equations $x'=Ax$ it is $\operatorname{Re}\lambda<0$,
and the proofs are short once we know how to measure the size of $A^k$. But eigenvalues do not tell the whole story: non-normal
matrices can grow enormously before they decay, which matters for every engineer and every numerical analyst. We then meet
Lyapunov's way of proving stability without computing eigenvalues, apply all of this to the speed of iterative methods (gradient
descent, Chebyshev's acceleration and Polyak's heavy ball, which together explain why ``momentum'' works), and end with Kalman's
criterion for when a system can be steered anywhere at all, and with Maxwell's cubic.

These stability estimates also underpin iterative methods and optimisation algorithms.

\begin{think}
The matrix $A=\bigl(\begin{smallmatrix}0.9&10\\0&0.9\end{smallmatrix}\bigr)$ has both eigenvalues equal to $0.9$, so $A^k\to0$. Roughly
how large does $\|A^k\|$ get before it starts to decrease? Compute $A^2$ and $A^3$ and guess a formula for $A^k$.
\end{think}

\section{Powers of a matrix and Gelfand's formula}

Recall that the \emph{spectral radius} of a square matrix is $\rho(A)=\max|\lambda|$ over its eigenvalues. For any vector norm on
$\C^n$ we use the induced operator norm $\|A\|=\max_{\|x\|=1}\|Ax\|$, which satisfies $\|AB\|\le\|A\|\|B\|$; if $Ax=\lambda x$ with
$x\ne0$, then $\|A\|\ge|\lambda|$, so $\rho(A)\le\|A\|$ for every induced norm. The key lemma says that this inequality is almost an
equality for a well-chosen norm.

\begin{lemma}[a norm adapted to $A$]\label{linear:lem:adaptednorm}
For every $A\in M_n(\C)$ and $\varepsilon>0$ there is an invertible matrix $S$ such that $\|S^{-1}AS\|_2\le\rho(A)+\varepsilon$.
Consequently the norm $\|x\|_S=\|S^{-1}x\|_2$ has induced norm $\|A\|_S\le\rho(A)+\varepsilon$.
\end{lemma}
\begin{proof}
By Schur's theorem (Theorem~\ref{linear:thm:schur}), $A=UTU^*$ with $U$ unitary and $T$ upper triangular, $T=\Lambda+N$ with
$\Lambda$ diagonal (the eigenvalues) and $N$ strictly upper triangular. For $\delta>0$ let $D=\operatorname{diag}(1,\delta,\delta^2,
\dots,\delta^{n-1})$. Then $D^{-1}TD$ has entries $t_{ij}\delta^{j-i}$: the diagonal is unchanged and every entry above it is
multiplied by at least $\delta$. So $D^{-1}TD=\Lambda+N_\delta$ with $\|N_\delta\|_2\le\|N_\delta\|_F\le\delta\|N\|_F$ for $\delta\le1$.
Hence $\|D^{-1}TD\|_2\le\|\Lambda\|_2+\delta\|N\|_F=\rho(A)+\delta\|N\|_F$, and $S=UD$ with $\delta$ small enough does it. Finally
$\|Ax\|_S=\|S^{-1}AS\,S^{-1}x\|_2\le(\rho(A)+\varepsilon)\|x\|_S$.
\end{proof}

\begin{theorem}[convergence of powers]\label{linear:thm:powers}
$A^k\to0$ as $k\to\infty$ if and only if $\rho(A)<1$. More precisely, for every $\varepsilon>0$ there is $C_\varepsilon$ with
$\|A^k\|\le C_\varepsilon(\rho(A)+\varepsilon)^k$ for all $k\ge0$.
\end{theorem}
\begin{proof}
With $S$ from Lemma~\ref{linear:lem:adaptednorm}, $A^k=S(S^{-1}AS)^kS^{-1}$, so $\|A^k\|_2\le\|S\|_2\|S^{-1}\|_2(\rho+\varepsilon)^k$. Every norm on
$M_n(\C)$ is bounded by a constant times $\|\cdot\|_2$, which gives the bound in any norm. If $\rho<1$, choose $\varepsilon$ with
$\rho+\varepsilon<1$. Conversely, if $Ax=\lambda x$ with $|\lambda|\ge1$ and $x\ne0$, then $\|A^kx\|=|\lambda|^k\|x\|\not\to0$.
\end{proof}

\begin{theorem}[Gelfand's formula, 1941]\label{linear:thm:gelfand}
For every norm on $M_n(\C)$, $\displaystyle\rho(A)=\lim_{k\to\infty}\|A^k\|^{1/k}$.
\end{theorem}
\begin{proof}
All norms on $M_n(\C)$ are equivalent, and $c^{1/k}\to1$ for every constant $c>0$, so it suffices to treat the operator norm
$\|\cdot\|_2$. Since $\rho(A^k)=\rho(A)^k$ and $\rho\le\|\cdot\|_2$, we have $\rho(A)\le\|A^k\|_2^{1/k}$. By Theorem~\ref{linear:thm:powers},
$\|A^k\|_2^{1/k}\le C_\varepsilon^{1/k}(\rho+\varepsilon)$, whose $\limsup$ is at most $\rho+\varepsilon$, for every $\varepsilon>0$.
\end{proof}

So the spectral radius is the \emph{asymptotic} growth rate of the powers. The opening ``Pause and think'' shows that it says nothing
about the \emph{transient}. For $A=\lambda I+cN$ with $N=\bigl(\begin{smallmatrix}0&1\\0&0\end{smallmatrix}\bigr)$, $N^2=0$ and the
binomial formula give
\[
A^k=\lambda^kI+k\lambda^{k-1}cN=\begin{pmatrix}\lambda^k&kc\lambda^{k-1}\\0&\lambda^k\end{pmatrix}.
\]
For $\lambda=0.9$, $c=10$, the corner $10k\cdot0.9^{k-1}$ reaches about $38.7$ at $k=9$ and $10$ before decaying, and $\|A^k\|$ stays above
$1$ for more than sixty steps (Figure~\ref{linear:fig:transient}). For normal matrices there is no such hump: $\|A^k\|_2=\rho(A)^k$ exactly,
since $A$ is unitarily diagonalisable. Transient growth is a phenomenon of non-normality, and it is not a curiosity: in fluid
dynamics, linearly stable flows (pipe flow, for instance) become turbulent at moderate speeds partly because tiny disturbances are
amplified by thousands before they decay, and the non-linear terms take over in between.

\begin{figure}[ht]
\centering
\fig{linear-ch12-transient}
\caption{$\|A^k\|_2$ for $A=\bigl(\begin{smallmatrix}0.9&10\\0&0.9\end{smallmatrix}\bigr)$ (blue, log scale) and for the normal matrix
$0.9I$ with the same eigenvalues (red). Both decay like $0.9^k$ in the end, but the non-normal one first grows by a factor of almost
$40$.}
\label{linear:fig:transient}
\end{figure}

\section{Differential equations and stability}

The solution of $x'=Ax$, $x(0)=x_0$ is $x(t)=e^{tA}x_0$ (Proposition~\ref{linear:prop:exp}); it is unique, since if $x'=Ax$ then
$\frac d{dt}(e^{-tA}x(t))=0$. If $Av=\lambda v$, then $e^{tA}v=e^{\lambda t}v$, so solutions starting on eigenvectors grow or decay
like $|e^{\lambda t}|=e^{t\operatorname{Re}\lambda}$. The \emph{spectral abscissa} $\alpha(A)=\max\operatorname{Re}\lambda$ plays the
role of $\log\rho$, and a matrix with $\alpha(A)<0$ is called \emph{stable} (or \emph{Hurwitz}).

\begin{theorem}[stability]\label{linear:thm:stability}
For every $\varepsilon>0$ there is $C_\varepsilon$ with $\|e^{tA}\|\le C_\varepsilon e^{(\alpha(A)+\varepsilon)t}$ for all $t\ge0$. In
particular every solution of $x'=Ax$ tends to $0$ if and only if $A$ is stable, and then it does so exponentially fast.
\end{theorem}
\begin{proof}
As in Lemma~\ref{linear:lem:adaptednorm}, write $S^{-1}AS=\Lambda+N_\delta$ with $\Lambda$ diagonal and $\|N_\delta\|_2\le\varepsilon$. Let
$y(t)=S^{-1}x(t)$, so $y'=(\Lambda+N_\delta)y$. Then
\[
\frac d{dt}\|y\|^2=2\operatorname{Re}\,y^*(\Lambda+N_\delta)y\le2\Bigl(\max_i\operatorname{Re}\lambda_i\Bigr)\|y\|^2+2\|N_\delta\|_2\|y\|^2\le
2(\alpha+\varepsilon)\|y\|^2.
\]
So $\frac d{dt}\bigl(e^{-2(\alpha+\varepsilon)t}\|y\|^2\bigr)\le0$, and $\|y(t)\|\le e^{(\alpha+\varepsilon)t}\|y(0)\|$. Returning to $x=Sy$,
$\|x(t)\|\le\|S\|\|S^{-1}\|e^{(\alpha+\varepsilon)t}\|x_0\|$. If $\alpha<0$, choose $\varepsilon<-\alpha$. Conversely, if $Av=\lambda v$ with
$\operatorname{Re}\lambda\ge0$, the solution $e^{\lambda t}v$ does not tend to $0$.
\end{proof}

For $2\times2$ real matrices the possible pictures are few, and Figure~\ref{linear:fig:phase} shows the main ones: real eigenvalues of the
same sign (a \emph{node}), complex eigenvalues (a \emph{spiral}, or a \emph{centre} if they are purely imaginary), and real eigenvalues
of opposite signs (a \emph{saddle}). Stability for $2\times2$ real matrices is $\operatorname{tr}A<0$ and $\det A>0$
(Exercise~\ref{linear:ex:hurwitz2}), the first case of the Routh--Hurwitz criterion.

\begin{figure}[ht]
\centering
\fig{linear-ch12-phase}
\caption{Phase portraits of $x'=Ax$ in the plane. Stable node and stable spiral: every trajectory tends to $0$. Centre: closed
orbits, neither stable nor unstable in the sense of Theorem~\ref{linear:thm:stability}. Saddle: trajectories come in along the stable
eigenvector and leave along the unstable one (red).}
\label{linear:fig:phase}
\end{figure}

\section{Lyapunov's method}

In his doctoral thesis of 1892, \emph{The General Problem of the Stability of Motion}, Aleksandr Lyapunov in Kharkov proposed to
prove stability without solving the equations: find an ``energy'' that decreases along every trajectory. For linear systems,
quadratic energies suffice, and the search for one is a linear equation.

\begin{theorem}[Lyapunov, 1892]\label{linear:thm:lyapunov}
Let $A\in M_n(\C)$. The following are equivalent.
\begin{enumerate}[label=(\roman*)]
\item $A$ is stable.
\item For every $Q\succ0$ there is $P\succ0$ with $A^*P+PA=-Q$.
\item There is $P\succ0$ with $A^*P+PA\prec0$.
\end{enumerate}
When $A$ is stable, the solution in (ii) is unique and equals $P=\int_0^\infty e^{tA^*}Qe^{tA}\,dt$.
\end{theorem}
\begin{proof}
(i)$\Rightarrow$(ii). By Theorem~\ref{linear:thm:stability} the integrand is bounded by $C^2\|Q\|e^{-2\eta t}$ with $\eta>0$, so the integral
converges. $P$ is Hermitian, and $x^*Px=\int_0^\infty(e^{tA}x)^*Q(e^{tA}x)\,dt>0$ for $x\ne0$. Since $\frac d{dt}\bigl(e^{tA^*}Qe^{tA}\bigr)=
A^*e^{tA^*}Qe^{tA}+e^{tA^*}Qe^{tA}A$,
\[
A^*P+PA=\int_0^\infty\frac d{dt}\bigl(e^{tA^*}Qe^{tA}\bigr)\,dt=\Bigl[e^{tA^*}Qe^{tA}\Bigr]_0^\infty=-Q.
\]
Uniqueness: the equation is $A^*P-P(-A)=-Q$, a Sylvester equation (Theorem~\ref{linear:thm:sylvester}); the eigenvalues of $A^*$ have negative
real parts and those of $-A$ positive ones, so they are disjoint.

(ii)$\Rightarrow$(iii) is clear. (iii)$\Rightarrow$(i). Let $Q=-(A^*P+PA)\succ0$ and $V(x)=x^*Px$. Along a solution,
\[
\frac d{dt}V(x(t))=x'^*Px+x^*Px'=x^*(A^*P+PA)x=-x^*Qx\le-\frac{\lambda_{\min}(Q)}{\lambda_{\max}(P)}V(x).
\]
So $V(x(t))\le e^{-ct}V(x_0)$ with $c>0$, and since $V(x)\ge\lambda_{\min}(P)\|x\|^2$, every solution tends to $0$: $A$ is stable by
Theorem~\ref{linear:thm:stability}.
\end{proof}

The ellipsoids $\{x^*Px\le r\}$ are traps: trajectories cross their boundaries inward and never come back. The matrix $P$ is computed
by solving a linear system, without any eigenvalue, which is how control engineers certify stability; and unlike eigenvalues, a
Lyapunov function survives small non-linear perturbations, which is how Lyapunov proved that the linearisation decides the stability
of a non-linear system near an equilibrium (when it is stable). The discrete-time version, for $x_{k+1}=Ax_k$, is the \emph{Stein
equation} $P-A^*PA=Q$ (Exercise~\ref{linear:ex:stein}).

\section{The speed of iterative methods}

To solve $Ax=b$ with $A$ symmetric positive definite, or equivalently to minimise $f(x)=\frac12x^{\mathsf T}Ax-b^{\mathsf T}x$, whose gradient
is $Ax-b$, one can avoid factorising $A$ and only multiply by it. This is how large sparse problems are solved, from structural
mechanics to the training of machine-learning models. Throughout, the eigenvalues of $A$ lie in $[m,L]$ with $0<m\le L$, and
$\kappa=L/m$ is the \emph{condition number}. Let $x_\star=A^{-1}b$ and $e_k=x_k-x_\star$.

\emph{Gradient descent} $x_{k+1}=x_k-\eta(Ax_k-b)$ gives $e_{k+1}=(I-\eta A)e_k$, a linear dynamical system. Its rate is
$\rho(I-\eta A)=\max_{\lambda\in\operatorname{spec}A}|1-\eta\lambda|$, and $\|e_k\|_2\le\rho^k\|e_0\|_2$ exactly, as $I-\eta A$ is symmetric.
The best step is $\eta=\frac2{m+L}$ (Exercise~\ref{linear:ex:gdstep}), with rate $\frac{L-m}{L+m}=\frac{\kappa-1}{\kappa+1}\approx1-\frac2\kappa$: the
number of iterations needed to gain a factor $10$ is proportional to $\kappa$.

Any method that builds $x_k$ from $x_0$ and $k$ multiplications by $A$, as linear combinations of $x_0,Ax_0-b,\dots$, has an error of the
form $e_k=p_k(A)e_0$, where $p_k$ is a polynomial of degree at most $k$ with $p_k(0)=1$. Then
\[
\|e_k\|_2\le\max_{\lambda\in[m,L]}|p_k(\lambda)|\,\|e_0\|_2,
\]
and the best such bound is a classical problem in approximation theory, solved by Chebyshev's polynomials $T_k(\cos\theta)=\cos k\theta$
(Volume~4 and Exercise~\ref{linear:ex:chebyshevrec}).

\begin{theorem}[Chebyshev acceleration]\label{linear:thm:chebyshevaccel}
Let $\ell(\lambda)=\frac{L+m-2\lambda}{L-m}$, which maps $[m,L]$ onto $[-1,1]$, and $x_0=\ell(0)=\frac{\kappa+1}{\kappa-1}$. The polynomial
$p_k(\lambda)=T_k(\ell(\lambda))/T_k(x_0)$ satisfies $p_k(0)=1$ and
\[
\max_{[m,L]}|p_k|=\frac1{T_k(x_0)}\le2\Bigl(\frac{\sqrt\kappa-1}{\sqrt\kappa+1}\Bigr)^k,
\]
and no polynomial $q$ of degree at most $k$ with $q(0)=1$ has a smaller maximum on $[m,L]$.
\end{theorem}
\begin{proof}
On $[-1,1]$, $|T_k|\le1$, so $\max_{[m,L]}|p_k|\le1/T_k(x_0)$, with equality since $T_k(1)=1$. For $x\ge1$, write $x=\cosh u$; then
$T_k(x)=\cosh ku$ (both sides satisfy the recurrence $T_{k+1}=2xT_k-T_{k-1}$ with the same initial values), so $T_k(x)\ge\frac12e^{ku}=
\frac12\bigl(x+\sqrt{x^2-1}\bigr)^k$. For $x=x_0$,
\[
x_0+\sqrt{x_0^2-1}=\frac{\kappa+1+2\sqrt\kappa}{\kappa-1}=\frac{(\sqrt\kappa+1)^2}{(\sqrt\kappa-1)(\sqrt\kappa+1)}=\frac{\sqrt\kappa+1}{\sqrt\kappa-1},
\]
which gives the bound. For optimality, $T_k(\cos\theta)=\cos k\theta$ takes the values $\pm1$ alternately at the $k+1$ points $\cos\frac{j\pi}
k$, $j=0,\dots,k$, so $p_k$ takes the values $\pm M$ alternately at $k+1$ points $\lambda_0<\dots<\lambda_k$ of $[m,L]$, where $M=1/T_k(x_0)$.
If $\max_{[m,L]}|q|<M$, then $r=p_k-q$ has the sign of $p_k$ at each $\lambda_j$, so it changes sign $k$ times and has $k$ roots in
$(m,L)$; it also vanishes at $0\notin[m,L]$. A polynomial of degree at most $k$ with $k+1$ roots is zero, so $q=p_k$, a contradiction.
\end{proof}

The number of iterations now scales like $\sqrt\kappa$ instead of $\kappa$: for $\kappa=10^4$, a hundred times fewer. The polynomials
$p_k$ satisfy a three-term recurrence (inherited from $T_k$), so the \emph{Chebyshev iteration} computes $x_{k+1}$ from $x_k$ and $x_{k-1}$
with one multiplication by $A$ per step. The conjugate gradient method of Hestenes and Stiefel (1952) does even better: it finds
the optimal polynomial for the actual spectrum of $A$ and the actual $e_0$ (in the norm $\|e\|_A=\sqrt{e^{\mathsf T}Ae}$), without knowing $m$
and $L$.

The simplest method with this square-root speed-up keeps a constant ``momentum''. \emph{Polyak's heavy-ball method} (1964) is
\[
x_{k+1}=x_k-\alpha(Ax_k-b)+\beta(x_k-x_{k-1}).
\]

\begin{theorem}[heavy ball on quadratics]\label{linear:thm:heavyball}
With $\alpha=\frac4{(\sqrt L+\sqrt m)^2}$ and $\beta=\bigl(\frac{\sqrt\kappa-1}{\sqrt\kappa+1}\bigr)^2$, the pairs $(e_{k+1},e_k)$ evolve by a
linear map whose spectral radius is $\sqrt\beta=\frac{\sqrt\kappa-1}{\sqrt\kappa+1}$. Hence $\|e_k\|\le C_\varepsilon(\sqrt\beta+\varepsilon)^k$ for every
$\varepsilon>0$.
\end{theorem}
\begin{proof}
The errors satisfy $e_{k+1}=\bigl((1+\beta)I-\alpha A\bigr)e_k-\beta e_{k-1}$. In an orthonormal eigenbasis of $A$ the coordinates decouple:
for an eigenvalue $\lambda$, the pair $(u_{k+1},u_k)$ is multiplied by $M_\lambda=\bigl(\begin{smallmatrix}1+\beta-\alpha\lambda&-\beta\\1&0
\end{smallmatrix}\bigr)$, whose characteristic polynomial is $z^2-(1+\beta-\alpha\lambda)z+\beta$. Put $s=\sqrt L$, $r=\sqrt m$, so that
$\sqrt\beta=\frac{s-r}{s+r}$. At $\lambda=m$, $1+\beta-\alpha m=\frac{2(s^2+r^2)-4r^2}{(s+r)^2}=2\sqrt\beta$; at $\lambda=L$, it equals $-2\sqrt\beta$;
being affine in $\lambda$, $|1+\beta-\alpha\lambda|\le2\sqrt\beta$ on $[m,L]$. So the discriminant is $\le0$: the roots are complex conjugate
(or double), with product $\beta$, hence both of modulus $\sqrt\beta$. The full linear map is block diagonal with blocks $M_\lambda$, so
its spectral radius is $\sqrt\beta$, and Theorem~\ref{linear:thm:powers} gives the bound.
\end{proof}

Figure~\ref{linear:fig:rates} compares the three methods for $\kappa=100$: gradient descent gains about $2\%$ per step, the heavy ball and the
Chebyshev bound about $18\%$. The heavy ball pays a transient price for its speed: its matrices $M_\lambda$ are not normal (at
$\lambda=m$ and $\lambda=L$ they are Jordan blocks), and the hump of Figure~\ref{linear:fig:transient} reappears. Nesterov's accelerated
gradient (1983) achieves the $\sqrt\kappa$ rate for all smooth strongly convex functions, not only quadratics; the momentum used to train
neural networks descends from these two ideas.

\begin{figure}[ht]
\centering
\fig{linear-ch12-rates}
\caption{Worst-case error over the spectrum $[1,100]$ ($\kappa=100$): gradient descent with the optimal step (rate $\frac{99}{101}$), the
heavy ball (rate $\frac9{11}$, with an initial hump), and the Chebyshev bound $2(\frac9{11})^k$.}
\label{linear:fig:rates}
\end{figure}

\section{Controllability}

So far the matrix has been left alone. In control theory one acts on it: $x'=Ax+Bu$, where $u(t)\in\C^p$ is an input we choose (thrust,
voltage, steam) and $B$ is $n\times p$. Starting from $x(0)=0$, the state at time $T$ is
\[
x(T)=\int_0^Te^{(T-s)A}Bu(s)\,ds,
\]
by the variation of constants formula (differentiate $e^{-tA}x(t)$). The system is \emph{controllable} if every state can be reached
in this way. In 1960 Rudolf Kalman, then at a research institute in Baltimore, gave a criterion that is pure linear algebra; the same
year he published the \emph{Kalman filter}, which a few years later was guiding the Apollo spacecraft to the Moon.

\begin{theorem}[Kalman's rank condition]\label{linear:thm:kalman}
For every $T>0$, the set of states reachable at time $T$ from $0$ (with continuous inputs $u$) is the column space of the
$n\times np$ matrix $K=\bigl(B\ \ AB\ \ A^2B\ \cdots\ A^{n-1}B\bigr)$. So the system is controllable if and only if $\operatorname{rank}K=n$.
\end{theorem}
\begin{proof}
Let $R_T$ be the reachable set, a subspace (the map $u\mapsto x(T)$ is linear). By Cayley--Hamilton (Theorem~\ref{linear:thm:cayleyhamilton}),
$A^n$ and hence every power $A^j$ is a linear combination of $I,A,\dots,A^{n-1}$, so $e^{\tau A}B=\sum_{j<n}c_j(\tau)A^jB$ has its columns in
$\operatorname{col}K$, and so does every $x(T)$: $R_T\subseteq\operatorname{col}K$.

Conversely, let $v\perp R_T$. Take $u(s)=B^*e^{(T-s)A^*}v$; then
\[
0=v^*x(T)=\int_0^T\bigl\|B^*e^{(T-s)A^*}v\bigr\|^2ds,
\]
so $v^*e^{\tau A}B=0$ for all $\tau\in[0,T]$. Differentiating $j$ times at $\tau=0$ gives $v^*A^jB=0$ for all $j$, that is, $v\perp\operatorname{col}K$. So
$R_T^\perp\subseteq(\operatorname{col}K)^\perp$, and $\operatorname{col}K\subseteq R_T$.
\end{proof}

The proof also produces a control: the \emph{controllability Gramian}
\[
W_T=\int_0^Te^{\tau A}BB^*e^{\tau A^*}d\tau
\]
 is invertible when
$\operatorname{rank}K=n$, and $u(s)=B^*e^{(T-s)A^*}W_T^{-1}x_T$ steers $0$ to $x_T$ with the least energy $\int_0^T\|u\|^2$ (Exercise~\ref{linear:ex:mininput}).
The time $T$ plays no role in \emph{whether} a state can be reached, only in the cost: going faster costs more.

\begin{example}[the double integrator]\label{linear:exa:doubleintegrator}
A cart on a rail with position $q$ and applied force $u$ obeys $q''=u$. With $x=(q,q')$, $A=\bigl(\begin{smallmatrix}0&1\\0&0\end{smallmatrix}
\bigr)$ and $B=\binom01$, so $K=\bigl(\begin{smallmatrix}0&1\\1&0\end{smallmatrix}\bigr)$ has rank $2$: one force controls both position and
velocity. By contrast, two identical carts pushed by the \emph{same} force ($A$ block diagonal with two copies, $B$ with the two copies of
$\binom01$) give a $K$ of rank $2<4$: their difference is invisible to the input and can never be corrected.
\end{example}

\section{A gem: Maxwell's cubic}

Back to the governor. Maxwell's third-order model led to a characteristic polynomial $s^3+as^2+bs+c$ with real coefficients, and the
question whether all its roots have negative real parts. His answer is short, and so is its proof.

\begin{theorem}[Maxwell, 1868]\label{linear:thm:maxwell}
All roots of $s^3+as^2+bs+c$ ($a,b,c$ real) have negative real parts if and only if $a>0$, $b>0$, $c>0$ and $ab>c$.
\end{theorem}
\begin{proof}
A real cubic has a real root $r$, and $s^3+as^2+bs+c=(s-r)(s^2+ps+q)$ with $p,q$ real. Comparing coefficients, $a=p-r$, $b=q-rp$,
$c=-rq$. A quadratic $s^2+ps+q$ has both roots in $\{\operatorname{Re}s<0\}$ if and only if $p>0$ and $q>0$ (real roots: their sum $-p$ is
negative and their product $q$ positive; complex roots: their real part is $-p/2$ and $q=|s|^2>0$). So the cubic is stable if and only if
$r<0$, $p>0$, $q>0$. Now compute
\[
ab-c=(p-r)(q-rp)+rq=p\bigl(q-rp+r^2\bigr)=p\,(b+r^2).
\]
If the cubic is stable, then $a=p-r>0$, $b=q-rp>0$, $c=-rq>0$ and $ab-c=p(b+r^2)>0$. Conversely, assume $a,b,c>0$ and $ab>c$. Then
$p(b+r^2)>0$ with $b+r^2>0$, so $p>0$. Next $c=-rq>0$, so $r\ne0$ and $q$ has the sign of $-r$. If $r>0$, then $q<0$ and $b=q-rp<0$,
a contradiction. So $r<0$ and $q>0$: the cubic is stable.
\end{proof}

For degree $n$, the Routh--Hurwitz criterion asks that the $n$ leading principal minors of a matrix built from the coefficients (the
\emph{Hurwitz matrix}) be positive; for $n=3$ these minors are $a$, $ab-c$ and $c(ab-c)$, and Maxwell's conditions come back. The
general proof is longer and uses the argument principle or Sturm sequences; Lyapunov's theorem gives yet another route.

\begin{beauty}
A governor hunting on a steam engine, a fluid that turns turbulent while every eigenvalue says it should be calm, an algorithm that
crawls on an ill-conditioned problem until it is given momentum, and a spacecraft steered by a rank condition: all are the same
object, a matrix raised to a power or put in an exponential. The eigenvalues decide where things end, the geometry of the
eigenvectors decides what happens on the way, and a well-chosen polynomial decides how fast we get there.
\end{beauty}

\section{Exercises}

\begin{exercise}[$\star$]\label{linear:ex:gelfandex}
For $A=\bigl(\begin{smallmatrix}1&1\\0&1\end{smallmatrix}\bigr)$ compute $\|A^k\|_\infty$ (maximum absolute row sum) and check Gelfand's
formula. Show that a matrix with $\rho(A)=1$ can have bounded or unbounded powers, and that the powers are bounded if $A$ is
diagonalisable with $\rho(A)\le1$.
\end{exercise}

\begin{exercise}[$\star$]\label{linear:ex:hurwitz2}
Show that a real $2\times2$ matrix is stable if and only if $\operatorname{tr}A<0$ and $\det A>0$.
\end{exercise}

\begin{exercise}[$\star\star$]\label{linear:ex:stein}
Show that $\rho(A)<1$ if and only if for some (equivalently every) $Q\succ0$ there is $P\succ0$ with $P-A^*PA=Q$, and that then
$P=\sum_{k\ge0}(A^*)^kQA^k$.
\end{exercise}

\begin{exercise}[$\star$]\label{linear:ex:oscillator}
For the damped oscillator $q''+2\zeta q'+q=0$ with $\zeta>0$, write $x=(q,q')$ and $x'=Ax$. Show that the energy $E=\frac12(q^2+q'^2)$ is
non-increasing but that $A^{\mathsf T}I+IA$ is only negative semidefinite. Find $P\succ0$ with $A^{\mathsf T}P+PA=-I$.
\end{exercise}

\begin{exercise}[$\star$]\label{linear:ex:gdstep}
Show that $\max_{\lambda\in[m,L]}|1-\eta\lambda|$ is minimised at $\eta=\frac2{m+L}$, with value $\frac{\kappa-1}{\kappa+1}$, and that the
popular choice $\eta=\frac1L$ gives the rate $1-\frac1\kappa$.
\end{exercise}

\begin{exercise}[$\star\star$]\label{linear:ex:chebyshevrec}
Show that $T_0=1$, $T_1=x$, $T_{k+1}=2xT_k-T_{k-1}$ defines polynomials with $T_k(\cos\theta)=\cos k\theta$ and $T_k(\cosh u)=\cosh ku$, and that
$T_k$ has leading coefficient $2^{k-1}$ for $k\ge1$.
\end{exercise}

\begin{exercise}[$\star\star$]\label{linear:ex:heavyjordan}
Show that the heavy-ball matrix $M_m$ of Theorem~\ref{linear:thm:heavyball} has the double eigenvalue $\sqrt\beta$ and is not diagonalisable,
and deduce that $\|M_m^k\|$ grows like $k\beta^{k/2}$.
\end{exercise}

\begin{exercise}[$\star$]\label{linear:ex:kalmanosc}
Is the harmonic oscillator $q''+q=u$ controllable? And the system $x'=-x+u$, $y'=-2y+u$ (two different first-order systems driven by
one input)? And $x'=-x+u$, $y'=-y+u$?
\end{exercise}

\begin{exercise}[$\star\star$]\label{linear:ex:mininput}
Assume $\operatorname{rank}K=n$. Show that $W_T\succ0$, that $u_\star(s)=B^*e^{(T-s)A^*}W_T^{-1}x_T$ steers $0$ to $x_T$ in time $T$, and that every
other input $u$ doing so has $\int_0^T\|u\|^2\ge\int_0^T\|u_\star\|^2=x_T^*W_T^{-1}x_T$.
\end{exercise}

\begin{exercise}[$\star$]\label{linear:ex:routh}
Is $s^3+s^2+2s+8$ stable? And $s^3+3s^2+3s+1$? Find all real $k$ for which $s^3+2s^2+s+k$ is stable.
\end{exercise}

\begin{exercise}[$\star\star$]\label{linear:ex:lognorm}
Let $\mu(A)=\lambda_{\max}\bigl(\frac{A+A^*}2\bigr)$. Show that $\|e^{tA}\|_2\le e^{t\mu(A)}$ for $t\ge0$, so $\mu(A)<0$ implies stability.
Show by an example that a stable matrix can have $\mu(A)>0$, and relate this to transient growth.
\end{exercise}

\begin{exercise}[$\star\star\star$]\label{linear:ex:leontief}
Show that if $A\ge0$ entrywise, $\rho(A)<1$ if and only if $I-A$ is invertible with $(I-A)^{-1}\ge0$, and that then $(I-A)^{-1}=\sum_kA^k$.
(This is the condition for Leontief's input--output model of an economy to be productive.)
\end{exercise}

\section{Hints}

\begin{hint}{linear:ex:gelfandex}
$A^k=\bigl(\begin{smallmatrix}1&k\\0&1\end{smallmatrix}\bigr)$. For the diagonalisable case, $\|A^k\|\le\|S\|\|S^{-1}\|$.
\end{hint}
\begin{hint}{linear:ex:hurwitz2}
The eigenvalues are the roots of $s^2-(\operatorname{tr}A)s+\det A$; use the quadratic case in the proof of Theorem~\ref{linear:thm:maxwell}.
\end{hint}
\begin{hint}{linear:ex:stein}
Imitate the proof of Theorem~\ref{linear:thm:lyapunov}, with sums instead of integrals and $V(x_{k+1})-V(x_k)$ instead of a derivative.
\end{hint}
\begin{hint}{linear:ex:oscillator}
$A=\bigl(\begin{smallmatrix}0&1\\-1&-2\zeta\end{smallmatrix}\bigr)$. Write $P=\bigl(\begin{smallmatrix}p&q\\q&r\end{smallmatrix}\bigr)$ and solve three
linear equations.
\end{hint}
\begin{hint}{linear:ex:gdstep}
$|1-\eta\lambda|$ is convex in $\lambda$, so its maximum on $[m,L]$ is at an endpoint; balance the two endpoints.
\end{hint}
\begin{hint}{linear:ex:chebyshevrec}
$\cos(k+1)\theta+\cos(k-1)\theta=2\cos\theta\cos k\theta$, and the same for $\cosh$.
\end{hint}
\begin{hint}{linear:ex:heavyjordan}
At $\lambda=m$ the characteristic polynomial is $(z-\sqrt\beta)^2$; a $2\times2$ matrix with a double eigenvalue is diagonalisable only if it
is scalar.
\end{hint}
\begin{hint}{linear:ex:kalmanosc}
Compute $K=(B\ AB)$ in each case.
\end{hint}
\begin{hint}{linear:ex:mininput}
For the minimality, show that $\int_0^T\langle u-u_\star,u_\star\rangle=0$.
\end{hint}
\begin{hint}{linear:ex:routh}
Theorem~\ref{linear:thm:maxwell}.
\end{hint}
\begin{hint}{linear:ex:lognorm}
Compute $\frac d{dt}\|x(t)\|^2$ as in the proof of Theorem~\ref{linear:thm:stability}.
\end{hint}
\begin{hint}{linear:ex:leontief}
For the converse, use a non-negative eigenvector of $A^{\mathsf T}$ for the eigenvalue $\rho(A)$ (Exercise~\ref{linear:ex:nonneg}).
\end{hint}

\section{Solutions}

\begin{solution}{linear:ex:gelfandex}
\idea Powers of a Jordan block grow polynomially, which is invisible to the $k$-th root.
\writeup By induction $A^k=\bigl(\begin{smallmatrix}1&k\\0&1\end{smallmatrix}\bigr)$, so $\|A^k\|_\infty=k+1$ and $(k+1)^{1/k}\to1=\rho(A)$. The
powers of $A$ are unbounded, while those of $I$ (also $\rho=1$) are bounded. If $A=SDS^{-1}$ with $D$ diagonal and $|d_i|\le1$, then
$\|A^k\|_2\le\|S\|_2\|D^k\|_2\|S^{-1}\|_2\le\|S\|_2\|S^{-1}\|_2$ for all $k$.
\end{solution}

\begin{solution}{linear:ex:hurwitz2}
\idea Stability of a $2\times2$ matrix is stability of a quadratic polynomial.
\writeup The characteristic polynomial is $s^2-(\operatorname{tr}A)s+\det A$. As shown in the proof of Theorem~\ref{linear:thm:maxwell}, a real
quadratic $s^2+ps+q$ has both roots in the open left half-plane if and only if $p>0$ and $q>0$; here $p=-\operatorname{tr}A$, $q=\det A$.
\end{solution}

\begin{solution}{linear:ex:stein}
\idea The discrete Lyapunov function decreases by $x^*Qx$ at each step.
\writeup If $\rho(A)<1$, Theorem~\ref{linear:thm:powers} makes $\sum_k(A^*)^kQA^k$ converge (terms bounded by $C^2\|Q\|(\rho+\varepsilon)^{2k}$). Its sum
$P$ is $\succeq Q\succ0$, and $A^*PA=\sum_{k\ge1}(A^*)^kQA^k=P-Q$. Conversely, if $P\succ0$ and $P-A^*PA=Q\succ0$, let $V(x)=x^*Px$ and $x_{k+1}=
Ax_k$: then $V(x_{k+1})=V(x_k)-x_k^*Qx_k\le\bigl(1-\frac{\lambda_{\min}(Q)}{\lambda_{\max}(P)}\bigr)V(x_k)$, a geometric decrease with ratio
in $[0,1)$ (it is $\ge0$ because $V\ge0$). So $x_k\to0$ for every $x_0$, $A^k\to0$, and $\rho(A)<1$ by Theorem~\ref{linear:thm:powers}.
\end{solution}

\begin{solution}{linear:ex:oscillator}
\idea Energy is the physicist's Lyapunov function, but it only decreases when the velocity is non-zero; a slightly tilted quadratic
form decreases strictly.
\writeup $E=\frac12x^{\mathsf T}x$ and $\frac{dE}{dt}=qq'+q'(-q-2\zeta q')=-2\zeta q'^2\le0$, which vanishes when $q'=0$: indeed $A^{\mathsf T}+A=
\bigl(\begin{smallmatrix}0&0\\0&-4\zeta\end{smallmatrix}\bigr)$ is only negative semidefinite. With $P=\bigl(\begin{smallmatrix}p&q\\q&r\end{smallmatrix}
\bigr)$, $A^{\mathsf T}P+PA=-I$ reads $-2q=-1$, $p-2\zeta q-r=0$, $2q-4\zeta r=-1$. So $q=\frac12$, $r=\frac1{2\zeta}$, $p=\zeta+\frac1{2\zeta}$, and
\[
P=\begin{pmatrix}\zeta+\frac1{2\zeta}&\frac12\\\frac12&\frac1{2\zeta}\end{pmatrix},
\]
which is positive definite: $p>0$ and $\det P=\frac12+\frac1{4\zeta^2}-\frac14>0$.
\end{solution}

\begin{solution}{linear:ex:gdstep}
\idea Balance the slowest and the fastest mode.
\writeup $g(\lambda)=|1-\eta\lambda|$ is convex, so $\max_{[m,L]}g=\max(|1-\eta m|,|1-\eta L|)$. As $\eta$ grows the first decreases (while
$\eta\le1/m$) and the second increases (once $\eta\ge1/L$); the maximum is smallest when $1-\eta m=\eta L-1$, i.e.\ $\eta=\frac2{m+L}$, with
value $1-\frac{2m}{m+L}=\frac{L-m}{L+m}=\frac{\kappa-1}{\kappa+1}$. For $\eta=\frac1L$, the maximum is $\max(1-\frac mL,0)=1-\frac1\kappa$.
\end{solution}

\begin{solution}{linear:ex:chebyshevrec}
\idea The recurrence is the product-to-sum formula for cosines.
\writeup Induction: if $T_j(\cos\theta)=\cos j\theta$ for $j\le k$, then $T_{k+1}(\cos\theta)=2\cos\theta\cos k\theta-\cos(k-1)\theta=\cos(k+1)\theta$, by
$\cos(k+1)\theta+\cos(k-1)\theta=2\cos\theta\cos k\theta$. The identity $\cosh(k+1)u+\cosh(k-1)u=2\cosh u\cosh ku$ gives the second claim in the same
way. The leading coefficient doubles at each step from $T_1=x$, so it is $2^{k-1}$.
\end{solution}

\begin{solution}{linear:ex:heavyjordan}
\idea A double eigenvalue without two independent eigenvectors means a Jordan block, and Jordan blocks grow polynomially.
\writeup At $\lambda=m$, $1+\beta-\alpha m=2\sqrt\beta$ (proof of Theorem~\ref{linear:thm:heavyball}), so the characteristic polynomial of $M_m$ is
$z^2-2\sqrt\beta z+\beta=(z-\sqrt\beta)^2$. If $M_m$ were diagonalisable it would be $\sqrt\beta I$, but its $(2,1)$ entry is $1$. So $M_m=
P(\sqrt\beta I+N)P^{-1}$ with $N\ne0$, $N^2=0$, and $M_m^k=P(\beta^{k/2}I+k\beta^{(k-1)/2}N)P^{-1}$, whose norm is between two positive
constants times $k\beta^{k/2}$ for large $k$.
\end{solution}

\begin{solution}{linear:ex:kalmanosc}
\idea Controllability fails exactly when some combination of the states does not feel the input.
\writeup Oscillator: $A=\bigl(\begin{smallmatrix}0&1\\-1&0\end{smallmatrix}\bigr)$, $B=\binom01$, $K=\bigl(\begin{smallmatrix}0&1\\1&0\end{smallmatrix}
\bigr)$, rank $2$: controllable. Second system: $A=\operatorname{diag}(-1,-2)$, $B=\binom11$, $AB=\binom{-1}{-2}$, $K$ has determinant $-2+1=-1\ne0$:
controllable, because the two modes respond differently to the same input. Third: $A=-I$, $B=\binom11$, $AB=-B$, rank $1$: not controllable;
$x-y$ satisfies $(x-y)'=-(x-y)$ whatever $u$ does.
\end{solution}

\begin{solution}{linear:ex:mininput}
\idea $u_\star$ is the input of minimal norm in an affine space of inputs, so it is orthogonal to the directions that do not change the
endpoint.
\writeup $v^*W_Tv=\int_0^T\|B^*e^{\tau A^*}v\|^2d\tau$, which vanishes only if $v^*e^{\tau A}B=0$ on $[0,T]$, hence (as in the proof of
Theorem~\ref{linear:thm:kalman}) only if $v\perp\operatorname{col}K=\C^n$: so $W_T\succ0$. With $\tau=T-s$, $\int_0^Te^{(T-s)A}Bu_\star(s)\,ds=W_TW_T^{-1}x_T=x_T$.
If $u$ also reaches $x_T$, then $d=u-u_\star$ satisfies $\int_0^Te^{(T-s)A}Bd(s)\,ds=0$, and
\[
\int_0^T\langle d,u_\star\rangle\,ds=\int_0^T\bigl\langle d(s),B^*e^{(T-s)A^*}w\bigr\rangle\,ds=\Bigl\langle\int_0^Te^{(T-s)A}Bd(s)\,ds,\,w\Bigr\rangle=0,
\]
with $w=W_T^{-1}x_T$. So $\int\|u\|^2=\int\|u_\star\|^2+\int\|d\|^2\ge\int\|u_\star\|^2$, and $\int\|u_\star\|^2=w^*W_Tw=x_T^*W_T^{-1}x_T$.
\end{solution}

\begin{solution}{linear:ex:routh}
\idea Three sign conditions and one product.
\writeup $s^3+s^2+2s+8$: $ab=2<8=c$, unstable (it has two roots with positive real part). $s^3+3s^2+3s+1=(s+1)^3$:
$ab=9>1$, stable. $s^3+2s^2+s+k$: we need $k>0$ and $2\cdot1>k$, so $0<k<2$. At $k=2$, $s^3+2s^2+s+2=(s+2)(s^2+1)$ has the roots $\pm i$: the
system oscillates forever, which is the onset of Maxwell's ``hunting''.
\end{solution}

\begin{solution}{linear:ex:lognorm}
\idea The symmetric part of $A$ measures the instantaneous growth rate of $\|x\|$; the eigenvalues only measure the long-run rate.
\writeup For $x'=Ax$, $\frac d{dt}\|x\|^2=2\operatorname{Re}x^*Ax=x^*(A+A^*)x\le2\mu(A)\|x\|^2$, so $\|x(t)\|\le e^{t\mu(A)}\|x(0)\|$ and $\|e^{tA}\|_2\le
e^{t\mu(A)}$. For $A=\bigl(\begin{smallmatrix}-1&10\\0&-1\end{smallmatrix}\bigr)$, stable, the symmetric part $\bigl(\begin{smallmatrix}-1&5\\5&-1\end{smallmatrix}
\bigr)$ has eigenvalue $4$, so $\mu(A)=4>0$: some solutions grow at first (initially at rate $e^{4t}$), the continuous-time version of
transient growth; for normal $A$, $\mu(A)=\alpha(A)$ and there is none.
\end{solution}

\begin{solution}{linear:ex:leontief}
\idea For non-negative matrices, the Neumann series has non-negative terms, and the Perron eigenvector turns ``positive inverse'' into
``spectral radius below one''.
\writeup If $\rho(A)<1$, then $\sum_kA^k$ converges (Theorem~\ref{linear:thm:powers}), $(I-A)\sum_{k\le K}A^k=I-A^{K+1}\to I$, so $(I-A)^{-1}=\sum_kA^k\ge0$.
Conversely, assume $(I-A)^{-1}\ge0$. By the Perron--Frobenius theorem for non-negative matrices (Exercise~\ref{linear:ex:nonneg}), $A^{\mathsf T}$ has an
eigenvector $y\ge0$, $y\ne0$, with $A^{\mathsf T}y=\rho y$, $\rho=\rho(A)$. Then $(I-A^{\mathsf T})y=(1-\rho)y$; as $I-A$ is invertible, $\rho\ne1$ and
$y=(1-\rho)(I-A^{\mathsf T})^{-1}y$. The vector $(I-A^{\mathsf T})^{-1}y=((I-A)^{-1})^{\mathsf T}y$ is $\ge0$ and non-zero, and $y\ge0$ is non-zero, so
$1-\rho>0$.
\end{solution}

\section*{Sources and further reading}

\begin{itemize}
\item J.~C.~Maxwell, \emph{On governors}, Proceedings of the Royal Society of London 16 (1868), 270--283; E.~J.~Routh, \emph{A Treatise on
the Stability of a Given State of Motion}, Macmillan, 1877; A.~Hurwitz, \emph{\"Uber die Bedingungen, unter welchen eine Gleichung nur
Wurzeln mit negativen reellen Theilen besitzt}, Mathematische Annalen 46 (1895), 273--284.
\item A.~M.~Lyapunov, \emph{The General Problem of the Stability of Motion} (1892), English translation, Taylor \& Francis, 1992.
\item I.~Gelfand, \emph{Normierte Ringe}, Matematicheskii Sbornik 9 (1941), 3--24.
\item L.~N.~Trefethen and M.~Embree, \emph{Spectra and Pseudospectra: The Behavior of Nonnormal Matrices and Operators}, Princeton
University Press, 2005; L.~N.~Trefethen, A.~E.~Trefethen, S.~C.~Reddy and T.~A.~Driscoll, \emph{Hydrodynamic stability without
eigenvalues}, Science 261 (1993), 578--584.
\item B.~T.~Polyak, \emph{Some methods of speeding up the convergence of iteration methods}, USSR Computational Mathematics and
Mathematical Physics 4 (1964), 1--17; Y.~Nesterov, \emph{A method of solving a convex programming problem with convergence rate
$O(1/k^2)$}, Soviet Mathematics Doklady 27 (1983), 372--376.
\item M.~R.~Hestenes and E.~Stiefel, \emph{Methods of conjugate gradients for solving linear systems}, Journal of Research of the National
Bureau of Standards 49 (1952), 409--436; Y.~Saad, \emph{Iterative Methods for Sparse Linear Systems}, 2nd ed., SIAM, 2003.
\item R.~E.~Kalman, \emph{On the general theory of control systems}, Proceedings of the First IFAC Congress, Moscow, 1960; R.~E.~Kalman,
\emph{A new approach to linear filtering and prediction problems}, Journal of Basic Engineering 82 (1960), 35--45; E.~D.~Sontag,
\emph{Mathematical Control Theory}, 2nd ed., Springer, 1998.
\item O.~Mayr, \emph{The Origins of Feedback Control}, MIT Press, 1970 (Watt's governor and its ancestors).
\end{itemize}
